\documentclass{article}
\usepackage[T1]{fontenc}
\usepackage{lmodern}
\usepackage{graphicx}
\usepackage{amsmath,amssymb}
\usepackage[a4paper,margin=2cm]{geometry}
\usepackage[absolute,overlay]{textpos}
\setlength{\TPHorizModule}{1mm}
\setlength{\TPVertModule}{1mm}
\usepackage[indonesian]{babel}
\usepackage{hyperref}
\setlength{\emergencystretch}{2em}
% unit: Halaman 49

\begin{document}

% === Halaman 49 (llm) ===

\begin{itemize}
    \item Kelebihan:
    \begin{itemize}
        \item Hanya perlu algoritma enkripsi saja
        \item Dapat mengenkripsi blok data secara acak (tidak harus sekuensial)
        \begin{itemize}
            \item Blok plainteks/cipherteks dapat diproses(enkripsi atau dekripsi) secara paralel
        \end{itemize}
        \item Sederhana, enkripsi dan dekripsi cepat
    \end{itemize}
    \vspace{10mm}
    \item Counter haruslah:
    \begin{itemize}
        \item Tidak diketahui atau tidak dapat diprediksi nilainya
    \end{itemize}
\end{itemize}

\end{document}
